\section{Objetivo del Informe}

El objetivo central del presente informe es evaluar los resultados y el impacto obtenidos tras la implementación del sistema ERP Odoo en Industria Textil Atlas E.I.R.L., analizando de manera cuantitativa y cualitativa la mejora en sus procesos comerciales, logísticos y de reabastecimiento. La intervención tecnológica tuvo como meta transformar el modelo operativo tradicional de la empresa a través de una plataforma integrada que centraliza el catálogo maestro de artículos deportivos con códigos SKU, agiliza la venta minorista de mostrador mediante terminales POS y optimiza el aprovisionamiento de stock sugerido por el sistema.

Asimismo, se buscó dotar a la organización de una infraestructura digital en la nube que reduzca drásticamente los tiempos de despacho y búsqueda en tienda (de 15--20 minutos a menos de 5 minutos), disminuya los errores por quiebre de inventario en un 20\%, y automatice la comunicación de órdenes de compra hacia los talleres confeccionistas externos. En este contexto, el informe presenta una evaluación de los beneficios reales alcanzados, el nivel de adopción por parte del personal de ventas y almacén, y el fortalecimiento de la capacidad gerencial para supervisar costos y fijar precios con márgenes de utilidad reales.
